\section{Введение и постановка задачи}

Рассматривается плоская задача трёх тел с равными массами \(m_1 = m_2 = m_3 = 1\) и гравитационной постоянной \(G = 1\). Исследуется однопараметрическое семейство периодических орбит (Ляпунова), ответвляющееся от эйлеровой коллинеарной конфигурации. В равновесии расстояния между телами равны 1, а система вращается с угловой скоростью \(\Omega = \sqrt{5/4}\).

Работа опирается на два набора данных, содержащих численно продолженные орбиты:
\begin{enumerate}
    \item \texttt{catalog\_unified.tsv} — 16\,128 орбит, разделённых на три блока (\textit{canonical}, \textit{near\_bif}, \textit{long}). Параметром продолжения в первых двух блоках служит начальное значение \(A = y_1(0)\) (в координатах Якоби), а в блоке \textit{long} — параметр в декартовых координатах. Состояния в этом каталоге уже повёрнуты на угол \(-\theta\).
    \item \texttt{catalog\_arclength.txt} — 83\,461 орбита. Параметром продолжения является псевдо-длина дуги в фазовом пространстве \((X, T)\). Состояния записаны в декартовых координатах без дополнительного поворота.
\end{enumerate}
Важно отметить, что с физической точки зрения каталоги описывают одно и то же семейство; различие в параметризации не вносит разрыва в физику орбит, а лишь меняет шаг продолжения.

Для интегрирования уравнений движения использовался высокоточный интегратор \texttt{heyoka} (метод Тейлора, тип данных \texttt{real128}, допуск \(\text{tol} = 10^{-25}\)). Для орбит с экстремально глубокими сближениями (\(d_{\min} < 10^{-7}\)) точность повышалась до 256 и более бит. Периодичность орбит уточнялась методом Ньютона для решения краевой задачи \(X(T; X_0) = X_0\).
